\def\tituloconcepto{QFT y estimación de fase}
% Si el capitulo se compila solo (sin main.tex), esta cabecera abre el documento.
% Antes de % Si el capitulo se compila solo (sin main.tex), esta cabecera abre el documento.
% Antes de % Si el capitulo se compila solo (sin main.tex), esta cabecera abre el documento.
% Antes de \input{cabecera} cada capitulo define \tituloconcepto.
\ifdefined\enmain\else
  \documentclass[11pt,a4paper]{article}
  \input{preambulo}
  \begin{document}
  \title{\tituloconcepto}
  \author{\autor}
  \date{\today}
  \maketitle
\fi
 cada capitulo define \tituloconcepto.
\ifdefined\enmain\else
  \documentclass[11pt,a4paper]{article}
  % Preambulo compartido por todos los documentos. Compilar con XeLaTeX/LuaLaTeX
% (o Tectonic). Si se usa Overleaf, elegir el compilador XeLaTeX.
\usepackage[spanish,es-noshorthands]{babel}
\usepackage{amsmath,amssymb}
\usepackage{booktabs,array}
\usepackage[a4paper,margin=2.5cm]{geometry}
\usepackage{xcolor}
\usepackage{listings}
\usepackage[most]{tcolorbox}
\usepackage{quantikz}
\usepackage[hidelinks]{hyperref}

% ---- Autor (editar a gusto) -------------------------------------------------
\newcommand{\autor}{Genesis Silesky}

% ---- Notacion de Dirac y atajos ---------------------------------------------
% quantikz ya define \ket y \bra; DeclareRobustCommand los redefine sin error.
\DeclareRobustCommand{\ket}[1]{\left|#1\right\rangle}
\DeclareRobustCommand{\bra}[1]{\left\langle#1\right|}
\DeclareRobustCommand{\braket}[2]{\left\langle#1\middle|#2\right\rangle}
\newcommand{\C}{\mathbb{C}}
\newcommand{\R}{\mathbb{R}}
\DeclareMathOperator{\Tr}{Tr}

% ---- Cajas de apoyo --------------------------------------------------------
% "concepto": definiciones previas que hacen falta para entender el capitulo.
\newtcolorbox{concepto}[1]{breakable,colback=blue!4,colframe=blue!55!black,
  fonttitle=\bfseries,title={Concepto previo: #1}}
% "nota": aclaraciones y advertencias.
\newtcolorbox{nota}[1]{breakable,colback=orange!6,colframe=orange!70!black,
  fonttitle=\bfseries,title={#1}}

% ---- Estilos de listados ----------------------------------------------------
\lstdefinestyle{py}{language=Python,basicstyle=\ttfamily\footnotesize,
  breaklines=true,postbreak=\mbox{\textcolor{gray}{$\hookrightarrow$}\space},
  numbers=left,numberstyle=\tiny\color{gray},frame=single,
  commentstyle=\color{green!40!black}\itshape,
  keywordstyle=\color{blue!70!black}\bfseries,
  stringstyle=\color{red!60!black},showstringspaces=false,tabsize=4,
  columns=fullflexible,keepspaces=true}
\lstdefinestyle{salida}{basicstyle=\ttfamily\footnotesize,breaklines=true,
  frame=leftline,backgroundcolor=\color{gray!8},columns=fullflexible,
  keepspaces=true}

  \begin{document}
  \title{\tituloconcepto}
  \author{\autor}
  \date{\today}
  \maketitle
\fi
 cada capitulo define \tituloconcepto.
\ifdefined\enmain\else
  \documentclass[11pt,a4paper]{article}
  % Preambulo compartido por todos los documentos. Compilar con XeLaTeX/LuaLaTeX
% (o Tectonic). Si se usa Overleaf, elegir el compilador XeLaTeX.
\usepackage[spanish,es-noshorthands]{babel}
\usepackage{amsmath,amssymb}
\usepackage{booktabs,array}
\usepackage[a4paper,margin=2.5cm]{geometry}
\usepackage{xcolor}
\usepackage{listings}
\usepackage[most]{tcolorbox}
\usepackage{quantikz}
\usepackage[hidelinks]{hyperref}

% ---- Autor (editar a gusto) -------------------------------------------------
\newcommand{\autor}{Genesis Silesky}

% ---- Notacion de Dirac y atajos ---------------------------------------------
% quantikz ya define \ket y \bra; DeclareRobustCommand los redefine sin error.
\DeclareRobustCommand{\ket}[1]{\left|#1\right\rangle}
\DeclareRobustCommand{\bra}[1]{\left\langle#1\right|}
\DeclareRobustCommand{\braket}[2]{\left\langle#1\middle|#2\right\rangle}
\newcommand{\C}{\mathbb{C}}
\newcommand{\R}{\mathbb{R}}
\DeclareMathOperator{\Tr}{Tr}

% ---- Cajas de apoyo --------------------------------------------------------
% "concepto": definiciones previas que hacen falta para entender el capitulo.
\newtcolorbox{concepto}[1]{breakable,colback=blue!4,colframe=blue!55!black,
  fonttitle=\bfseries,title={Concepto previo: #1}}
% "nota": aclaraciones y advertencias.
\newtcolorbox{nota}[1]{breakable,colback=orange!6,colframe=orange!70!black,
  fonttitle=\bfseries,title={#1}}

% ---- Estilos de listados ----------------------------------------------------
\lstdefinestyle{py}{language=Python,basicstyle=\ttfamily\footnotesize,
  breaklines=true,postbreak=\mbox{\textcolor{gray}{$\hookrightarrow$}\space},
  numbers=left,numberstyle=\tiny\color{gray},frame=single,
  commentstyle=\color{green!40!black}\itshape,
  keywordstyle=\color{blue!70!black}\bfseries,
  stringstyle=\color{red!60!black},showstringspaces=false,tabsize=4,
  columns=fullflexible,keepspaces=true}
\lstdefinestyle{salida}{basicstyle=\ttfamily\footnotesize,breaklines=true,
  frame=leftline,backgroundcolor=\color{gray!8},columns=fullflexible,
  keepspaces=true}

  \begin{document}
  \title{\tituloconcepto}
  \author{\autor}
  \date{\today}
  \maketitle
\fi


\section{Transformada de Fourier cuántica y estimación de fase}

\subsection{Objetivo}
Construir la transformada de Fourier cuántica (QFT) con un circuito y usarla
para estimar la fase de un autovalor (estimación de fase, QPE).

\subsection{Conceptos previos}
\begin{concepto}{fases, puertas controladas y números binarios}
Ver capítulos~05a (fases e interferencia) y~14 (puertas controladas).
Una fase se puede escribir como una fracción binaria:
$0.\varphi_1\varphi_2\ldots\varphi_m=\varphi_1/2+\varphi_2/4+\cdots+\varphi_m/2^m$.
Por ejemplo $0.75=0.11_2$ y $1/4=0.01_2$.
\end{concepto}

\subsection{Fundamento teórico}
\paragraph{QFT.} La QFT es un componente central de muchos algoritmos (Shor y la
estimación de fase, entre ellos). Es la versión cuántica de la transformada de
Fourier discreta (DFT). Con $N=2^n$ estados de la base, lleva
\begin{equation}
  \ket j\ \longmapsto\ \frac1{\sqrt N}\sum_{k=0}^{N-1}e^{2\pi i\,jk/N}\ket k ,
\end{equation}
una matriz unitaria con entradas $U[k,j]=e^{2\pi ijk/N}/\sqrt N$. Aunque la
fórmula se parece a la de la DFT, el significado es distinto: la DFT transforma
una función discreta en el dominio del tiempo y la QFT transforma un estado
cuántico en otro.

\paragraph{El circuito.} Escribiendo $j=j_1j_2\ldots j_n$ en binario, la QFT se
puede reescribir como un producto de qubits:
\begin{equation}
  \ket j\mapsto\frac1{2^{n/2}}\big(\ket0+e^{2\pi i\,0.j_n}\ket1\big)
  \big(\ket0+e^{2\pi i\,0.j_{n-1}j_n}\ket1\big)\cdots
  \big(\ket0+e^{2\pi i\,0.j_1j_2\ldots j_n}\ket1\big).
\end{equation}
Esto da un circuito con tres tipos de pieza:
\begin{enumerate}
  \item una Hadamard en cada qubit;
  \item fases controladas $CR_k$, que aplican la fase $e^{2\pi i/2^k}$ a $\ket1$
    ($CR_2$ es la puerta $S$ controlada) y trasladan la información de los
    demás bits a la fase;
  \item al final, puertas SWAP que invierten el orden de los qubits.
\end{enumerate}
El número de puertas crece como $O(n^2)$, es decir, $O((\log N)^2)$ para $N$
estados.

\paragraph{QPE.} Dada una unitaria $U$ y un autovector $\ket\psi$ con
$U\ket\psi=e^{2\pi i\varphi}\ket\psi$, la estimación de fase calcula $\varphi$.
La QFT pasa información del estado a la fase; la QPE pasa información de la
fase al estado. Usa dos registros: el primero con $m$ qubits (da una precisión de
$m$ bits) y el segundo, inicializado en el autovector $\ket\psi$. Pasos:
\begin{enumerate}
  \item Primer registro en $\ket0^{\otimes m}$ y segundo en $\ket\psi$.
  \item Hadamard en cada qubit del primer registro.
  \item Puertas $U^{2^k}$ controladas por el qubit $k$ del primer registro.
    Cada una deja una fase $e^{2\pi i\varphi2^k}$ en $\ket1$.
  \item QFT inversa sobre el primer registro.
  \item Se mide el primer registro: sale el número binario $\varphi_1\varphi_2\ldots\varphi_m$.
\end{enumerate}
Por ejemplo, la puerta $S$ tiene autovector $\ket1$ con autovalor
$i=e^{2\pi i\cdot1/4}$, así que $\varphi=1/4=0.01_2$.

\subsection{Código (completamente comentado)}
\lstinputlisting[style=py]{../../fuentes/17_qft_qpe.py}

\subsection{Explicación celda por celda}
\begin{description}
  \item[Celda 2] Construye la matriz de la QFT y comprueba que es unitaria.
  \item[Celda 3] Construye el circuito (Hadamard, fases controladas y SWAP) y
    comprueba que su matriz es igual a la de la QFT para $n=2,3,4$.
  \item[Celda 4] Aplica la QFT de dos qubits a cada estado de la base y muestra
    las fases de las cuatro amplitudes.
  \item[Celda 5] Define la QPE: Hadamard, fases controladas con
    $2\pi\varphi2^k$, QFT inversa y medición.
  \item[Celda 6] Prueba fases exactas: puertas $S$ ($1/4$), $T$ ($1/8$) y $Z$
    ($1/2$).
  \item[Celda 7] Prueba $\varphi=1/3$, que no se escribe exactamente con pocos
    bits, con $m=2,4,6$.
\end{description}

\subsection{Salida obtenida}
\lstinputlisting[style=salida]{../../fuentes/salidas/17_qft_qpe.txt}
(El dibujo del circuito no se reproduce aquí.) Las cuentas de la QPE dependen
del azar de las mediciones; el resultado más probable es estable.

\subsection{Interpretación}
El circuito coincide con la matriz de la QFT. Con 2 qubits, cada entrada da
cuatro amplitudes del mismo módulo y fases que son múltiplos de $90^\circ$; por
ejemplo $\ket{01}\to[0,90,180,270]$. Con las fases exactas, la QPE da el
resultado correcto con probabilidad~1: $S\to01$ ($0.25$), $T\to001$ ($0.125$)
y $Z\to10$ ($0.5$). Para $\varphi=1/3$, el resultado más probable se acerca
cada vez más: el error baja de $0.083$ ($m=2$) a $0.021$ ($m=4$) y a $0.005$
($m=6$), con probabilidad cercana a $0.69$.

\subsection{Ejercicios}
\begin{enumerate}
  \item Comprueba la QPE con $\varphi=3/8$ y $m=3$. ¿Qué cadena esperas?
  \item ¿Qué cadena da la QPE para la puerta $Z$ con $m=3$ bits? Predícelo
    y compruébalo.
  \item Cuenta las puertas del circuito de la QFT para $n=5$ y compara con
    $n^2$.
\end{enumerate}

\subsection{Errores comunes}
\begin{itemize}
  \item Olvidar los SWAP finales: sin ellos el circuito da la QFT con los
    qubits en orden inverso.
  \item Esperar un resultado exacto cuando la fase no cabe en $m$ bits:
    la QPE da el valor más cercano, con alta probabilidad pero no exacta.
\end{itemize}

\ifdefined\enmain\else
  \end{document}
\fi

